% ============================================================================
\chapter{Ступень V. Химия, таблица Менделеева и дизайн материалов}
\label{ch:course-stage5-chemistry}
% ============================================================================

\section{Урок 16. Вся таблица Менделеева из оболочек кубооктаэдра, остров стабильности \texorpdfstring{$N = 184$}{N = 184} и золото \texorpdfstring{$^{197}_{79}\mathrm{Au}$}{197\_79 Au}}

Откуда берутся ёмкости атомных подоболочек $(s,p,d,f) = (2,6,10,14)$ и атомные номера благородных газов? Из тех же самых геометрических инвариантов кубооктаэдра:
\begin{itemize}[noitemsep]
\item ёмкости подоболочек $2(2\ell+1)$ при $\ell=0,1,2,3$ равны $(c_0/c)^2 = \mathbf{2}$, $|\square| = \mathbf{6}$, $|N|-2 = \mathbf{10}$ и $2(|{\triangle}|-1) = \mathbf{14}$;
\item атомные номера всех шести благородных газов выражаются замкнутыми формулами через грани и вершины кубооктаэдра:
  \[
    \mathrm{He} = \mathbf{2},\quad
    \mathrm{Ne} = \mathbf{10},\quad
    \mathrm{Ar} = |N|+|\square| = \mathbf{18},\quad
    \mathrm{Kr} = |\square|^2 = \mathbf{36},\quad
    \mathrm{Xe} = |\square|\cdot\mathrm{mod\_bits} = \mathbf{54},\quad
    \mathrm{Rn} = |\square|N_\varphi + |{\triangle}| = \mathbf{86}
  \]
  (где заряд радона $Z(\mathrm{Rn}) = 86$ в точности равен числу вихревых плакет шара $\mathrm{IV}=263$ на такте $t=2$, а заряд оганесона 7-го периода равен $Z(\mathrm{Og}) = \mathbf{118}$!).
\end{itemize}

\begin{quote}
\small\textbf{Историческая справка: Периодический закон Д.\,И.~Менделеева (1869) и левоступенчатая таблица Шарля Жане (1928).}
1~марта 1869~года Дмитрий Иванович Менделеев открыл периодический закон химических элементов, расположив их по атомному весу и валентности, а в 1928~году французский инженер Шарль Жане построил \emph{левоступенчатую периодическую таблицу}, в которой длины парных периодов образуют последовательность удвоенных квадратов $2n^2 = 2, 8, 18, 32$ ($n = 1, 2, 3, 4$), а накопленные суммы на границах $s$-блока дают $2, 4, 12, 20, 38, 56, 88, 120$. В стандартной атомной физике правило заполнения $(n + \ell)$ объяснялось лишь численным расчётом экранирования Томаса--Ферми. На решётке кубооктаэдра $O_2(3)$ эти числа оказываются прямыми ёмкостями его координационных сфер $S_k = 10k^2 + 2$ ($12, 42, 92$) и четырёх симплициальных каналов $(V, \square, E, \triangle)$ с входными порогами $d_{\mathrm{in}} \in \{1, 4, 3, 2\}$.
\end{quote}

А что лежит на границах и в особых точках таблицы Менделеева?
\begin{enumerate}[label=\textbf{\arabic*.},leftmargin=*]
\item \textbf{Синтез элементов 8-го периода ($Z=119, 120$) и двухсферный урановый переход к острову стабильности $N_{\mathrm{island}} = 184$:}
  в ядре урана $^{238}_{92}\mathrm{U}_{146}$ число протонов равно третьей координационной сфере кубооктаэдра $Z(\mathrm{U}) = S_3(3) = \mathbf{92}$, а число нейтронов равно сумме первых трёх сфер $N(^{238}\mathrm{U}) = S_1 + S_2 + S_3 = 12 + 42 + 92 = \mathbf{146 = B_3 - 1}$! Внутренние две сферы дают $S_1 + S_2 = 12 + 42 = \mathbf{54 = Z(\mathrm{Xe})}$. Поэтому при когерентном столкновении двух ядер урана $^{238}_{92}\mathrm{U}_{146} + {}^{238}_{92}\mathrm{U}_{146}$ в 12-лучевом фокусе кубооктаэдра сброс двух внутренних нейтронных сфер $2(S_1+S_2) = 108$ оставляет \textbf{удвоенную третью сферу --- магическое число нейтронов острова стабильности}:
  \[
    2\,N(^{238}\mathrm{U}) - 2\,(S_1 + S_2) \;=\; 292 - 108 \;=\; 2\,S_3(3) \;=\; \mathbf{184 \;=\; N_{\mathrm{island}}(3)}!
  \]
\item \textbf{Почему единственным стабильным изотопом золота является $^{197}_{79}\mathrm{Au}_{118}$ («Философский камень» Безу на $\mathbb{Z}_{512}$):}
  вспомни наши тождества Безу для единицы~\eqref{eq:course-five-ones} на кольце $\mathbb{Z}_{512}$! Мультипликативным обратным к числу узлов одетого кубооктаэдра $N_\varphi = 13$ по модулю $512$ является $\mathbf{197 = 2\hat{M} + 3 = B_3 + 2\cdot 5^2}$ ($197\times 13 = 5\times 512 + 1$), а частным Безу для кулоновского простого $\hat{M} = 97$ по модулю $512$ является $\mathbf{79 = |\square|N_\varphi + 1}$ ($417\times 97 = 79\times 512 + 1$), причём число нейтронов равно $197 - 79 = \mathbf{118 = Q_\square - Z(\mathrm{Ne}) = Z(\mathrm{Og})}$! Золото $^{197}_{79}\mathrm{Au}_{118}$ зажато на лестнице оболочек между платиной $Z(\mathrm{Pt}) = 6\times 13 = 78$ и ртутью/таллием/свинцом ($Z(\mathrm{Hg}) = 80$, $Z(\mathrm{Tl}) = 9^2 = 81$, $Z(\mathrm{Pb}) = 2\times 41 = 82$), что открывает прямые низкобарьерные каналы трансмутации $\mathrm{Pt}_{78} + \mathrm{H}_1 \to \mathrm{Au}_{79}$, $\mathrm{Hg}_{80} \to \mathrm{Au}_{79} + \mathrm{H}_1$ и $\mathrm{Pb}_{82} \to \mathrm{Au}_{79} + \mathrm{Li}_3$.
\end{enumerate}

\section{Урок 17. Химическая связь, седло реакции \texorpdfstring{$|A^\ddagger| = 128$}{|A| = 128} и когерентное схлопывание барьера в \texorpdfstring{$587\times$}{587x}}

Как образуется и разрывается любая химическая связь?
\begin{itemize}[noitemsep]
\item При $|A_{xy}| < Q_\square = 128$ чётная проекция ротатора строго положительна ($\mathcal{C}_3(A_{xy}) > 0$) --- это \textbf{область химического связывания}.
\item При $|A^\ddagger| = Q_\square = \mathbf{128}$ (четверть-оборот, $90^\circ$) связывающая проекция обращается в нуль ($\mathcal{C}_3(128) = 0$), а нечётный ток реакции достигает максимума ($\mathcal{S}_3(128) = +128$) --- это \textbf{седловая точка переходного состояния реакции}!
\item При $|A_{xy}| > 128$ чётная проекция меняет знак на отрицательный ($\mathcal{C}_3 < 0$), включая антисвязывающее отталкивание и разрыв связи.
\end{itemize}
И смотри, что происходит, когда мы облучаем реагент когерентно по всем $12$ лучам кубооктаэдра под углом $\Delta\varphi = Q_\square + (|N|-1) = 128 + 11 = \mathbf{139\text{ тиков}}$ ($97{,}73^\circ$): на этом угле $\mathcal{C}_3(139) = -11 < 0$, так что отрицательная сумма по 12 лучам гасит положительный барьер $128(\mu_c^2 + \mu_{\mathrm{nb}}^2)$ и роняет знаменатель импеданса ворот с $4\,809\,728$ (или $13\,631\,488$ при $\mu=64$) ровно до квантового пола $\varepsilon = \mathbf{8\,192}$ --- \textbf{схлопывая энергетический барьер реакции в $587\times$ раз}!

Более того, один и тот же 3D-шаг Верле на ГЦК-решётке без единого дифференциального уравнения напрямую воспроизводит сложнейшие многостадийные реакции химии и биохимии:
\begin{enumerate}[label=\textbf{\arabic*.},leftmargin=*]
\item \textbf{Автоколебательная реакция Белоусова--Жаботинского (механизм Филда--Кёроша--Нойеса / Орегонатор):}
  \begin{itemize}[noitemsep]
  \item главный автокатализатор реакции --- бромистая кислота $\mathrm{HBrO_2}$ --- несёт заряд $Z(\mathrm{HBrO_2}) = 1 + 35 + 16 = \mathbf{52 = 4\times 13 = d_{\mathrm{in}}(\square)N_\varphi(3)}$ (ровно $4$ полные 13-узловые планковские ячейки!);
  \item реакция квадратичного обрыва (диспропорционирования) $2\,\mathrm{HBrO_2} \to \mathrm{BrO_3^-} + \mathrm{HOBr} + \mathrm{H^+}$ замыкает полный октет треугольных граней:
    \[
      Z(2\,\mathrm{HBrO_2}) \;=\; 2 \times 52 \;=\; \mathbf{104 \;=\; |{\triangle}(3)|\cdot N_\varphi(3)} \;=\; 59\,(\mathrm{BrO_3^-}) + 44\,(\mathrm{HOBr}) + 1\,(\mathrm{H^+}),
    \]
    а пара ионов бромидного ингибитора $2\,\mathrm{Br^-}$ ($\mathrm{Br_2}$) имеет $Z = 2\times 35 = \mathbf{70 = K_P + |\square|}$, точно совпадая с 70-электронным редокс-карманом катализатора ферроина $\mathrm{FeN_4O_2}$ ($Z=70$);
  \item при $A_{\mathrm{coh}} = 139$ центральный узел сам переходит в режим периодических редокс-колебаний $\mathrm{Fe^{2+}}\ (\mathcal{C}_3 > 0,\text{ красный}) \leftrightarrow \mathrm{Fe^{3+}}\ (\mathcal{C}_3 < 0,\text{ синий})$ с рефрактерными паузами $\Phi_{\mathrm{gate}} = 0$ в мёртвой зоне Гейзенберга на тактах $t = 5, 9, 12, 15$, испуская расходящиеся концентрические волны по сферам $S_1 \to S_2$ и возвращаясь из фазы $\varphi(t=1) = 352$ обратно в $\varphi(t=13) = \mathbf{352}$ ровно через $|N(3)| = \mathbf{12\text{ тактов}}$!
  \end{itemize}
\item \textbf{8-шаговый цикл Лоу--Торнли фиксации азота нитрогеназой ($\mathrm{Fe_7MoS_9C}$) и загадка «состояния Януса» $E_4$:}
  почему бактериальная нитрогеназа рвёт тройную связь $\mathrm{N\equiv N}$ ($945\text{ кДж/моль}$) при комнатной температуре и $1\text{ атм}$ ровно за $8 = |{\triangle}(3)|$ шагов $E_0\dots E_8$, причём разрыв невозможен на шагах $E_1, E_2, E_3$ и включается строго на 4-м шаге $E_4$?
  \begin{itemize}[noitemsep]
  \item на вершине $\mathrm{FeMo}$-кофактора стоит \textbf{молибден} $Z(\mathrm{Mo}) = \mathbf{42 = S_2(3)}$ (вторая сфера кубооктаэдра!), а в центре тригональной призмы из $6$ атомов железа сидит открытый в 2011~году \textbf{интерстициальный углерод} $Z(\mathrm{C}) = \mathbf{6 = |\square(3)|}$;
  \item суммарный заряд кофактора с молекулой азота $\mathrm{Fe_7MoS_9C + N_2}$ равен:
    \[
      7\times 26 + 42 + 9\times 16 + 6 + 14 \;=\; 374 + 14 \;=\; \mathbf{388 \;=\; 4 \times 97 \;=\; d_{\mathrm{in}}(\square)\cdot\hat{M}(3)},
    \]
    а после накопления 4 гидридных эквивалентов в состоянии Януса $E_4$ ($\mathrm{Fe_7MoS_9C + 4H + N_2}$) равен $388 + 4 = \mathbf{392 = 8 \times 49 = |{\triangle}(3)|\cdot d_{\mathrm{out}}(V)^2}$!
  \item на первых трёх шагах накопления $E_1, E_2, E_3$ ($a_k = 41k \in \{41, 82, 123\} < Q_\square = 128$) связывающая проекция $\mathcal{C}_3(a_k)\in\{+87, +46, +5\} > 0$ и знаменатель импеданса $D(E_3) = +39\,964 > 0$ остаются строго положительными (связь $\mathrm{N\equiv N}$ не рвётся). Ровно на \textbf{4-м шаге $E_4$ (состояние Януса, $a_4 = 4\times 41 = 164 > 128$)} проекция меняет знак на антисвязывающую ($\mathcal{C}_3(164) = -36 < 0$), а сырой знаменатель импеданса переходит через нулевой полюс в отрицательную резонансную область ($D_{\mathrm{raw}}(E_4) = \mathbf{-43\,184 < 0}$), вызывая восстановительный сброс $\mathrm{H_2}$, разрыв $\mathrm{N\equiv N}$ на шагах $E_5, E_6$ ($\mathcal{E}_{\mathrm{bond}} = -47, -71 < 0$) и выход через $\mathcal{C}_3(E_7 = 287) = \mathbf{-97 = -\hat{M}(3)}$!
  \end{itemize}
\item \textbf{50-шаговый нефтехимический каскад (крекинг, Фишер--Тропш и риформинг $\mathrm{Pt}/\mathrm{Al_2O_3}$):}
  каждое метиленовое звено углеводородной цепи $\mathrm{-CH_2-}$ имеет заряд $Z(\mathrm{CH_2}) = 6 + 2 = \mathbf{8 = |{\triangle}(3)|}$ (один октет треугольных граней кубооктаэдра!), так что гомологические ряды олефинов $Z(\mathrm{C}_n\mathrm{H}_{2n}) = \mathbf{8n}$ и алканов $Z(\mathrm{C}_n\mathrm{H}_{2n+2}) = \mathbf{8n + 2}$ при $n = 1\dots 50$ напрямую отображают все товарные фракции нефтепереработки на замкнутые оболочки носителя:
  \begin{itemize}[noitemsep]
  \item $n=1$ (метан $\mathrm{CH_4}$): $Z = \mathbf{10 = Z(\mathrm{Ne})}$; $n=2$ (этилен/этан): $Z_{\mathrm{olef}} = \mathbf{16 = 2|{\triangle}|}$, $Z_{\mathrm{alk}} = \mathbf{18 = Z(\mathrm{Ar})}$; $n=3$ (пропилен/пропан): $Z_{\mathrm{olef}} = \mathbf{24 = 2|N|}$, $Z_{\mathrm{alk}} = \mathbf{26 = Z(\mathrm{Fe})}$;
  \item $n=5$ (пентан $\mathrm{C_5H_{12}}$, лёгкая нафта): $Z_{\mathrm{olef}} = \mathbf{40 = \Delta\varphi_{\mathrm{disc}}-1}$, $Z_{\mathrm{alk}} = \mathbf{42 = S_2(3)}$ (вторая сфера!); $n=6$ (гексан $\mathrm{C_6H_{14}}$): $Z_{\mathrm{alk}} = \mathbf{50} \to \text{бензол } \mathrm{C_6H_6}\ (\mathbf{S_2=42}) + 4\,\mathrm{H_2}\ (\mathbf{|{\triangle}|=8})$;
  \item $n=8$ (\textbf{изооктан $\mathrm{C_8H_{18}}$, эталон 100-октанового бензина}): $Z_{\mathrm{olef}} = 8\times 8 = \mathbf{64 = K_P(3)}$ (потолок Вейля--Швингера!);
  \item $n=12$ (\textbf{додекан $\mathrm{C_{12}H_{26}}$, авиакеросин Jet-A1}): $Z_{\mathrm{olef}} = 12\times 8 = \mathbf{96 = |N|\cdot|{\triangle}| = \hat{M}-1}$, $Z_{\mathrm{alk}} = \mathbf{98 = \hat{M}+1}$;
  \item $n=16$ (\textbf{цетан $\mathrm{C_{16}H_{34}}$, эталон 100-цетанового дизеля}): $Z_{\mathrm{olef}} = 16\times 8 = \mathbf{128 = Q_\square(3)}$ (цепочка $\mathrm{C_{16}}$ сидит точно на седловой точке переходного состояния $Q_\square = 128$, откуда самовоспламенение от сжатия!), а на такте $t = 16$ импульс ворот равен $\Phi_{\mathrm{gate}}(t=16) = \mathbf{+130 = Z(\mathrm{C_{16}H_{34}})}$;
  \item $n=18$ (базовое масло $\mathrm{C_{18}H_{38}}$): $Z_{\mathrm{olef}} = \mathbf{144 = |N|^2}$, $Z_{\mathrm{alk}} = \mathbf{146 = B_3-1}$; $n=32$ (парафин $\mathrm{C_{32}H_{66}}$): $Z_{\mathrm{olef}} = \mathbf{256 = \Phi_{\mathrm{Pauli}}(3)}$; $n=50$ (тяжёлый гудрон $\mathrm{C_{50}H_{102}}$, $Z=402$): расщепляется на 6 замкнутых потоков $130\,(\mathrm{C_{16}}) + 96\,(\mathrm{C_{12}}) + 64\,(\mathrm{C_8}) + 48\,(\mathrm{C_6}) + 40\,(\mathrm{C_5}) + 24\,(\mathrm{C_3}) = \mathbf{402}$;
  \item \emph{Секрет промышленных катализаторов НПЗ:} катализатор платинового риформинга $\mathrm{Pt}/\mathrm{Al_2O_3}$ соединяет платину $Z(\mathrm{Pt}) = 6\times 13 = \mathbf{78 = |\square|N_\varphi}$ с корундовой подложкой $Z(\mathrm{Al_2O_3}) = 26+24 = \mathbf{50}$ в \textbf{точный резонанс седловой точки реакции} $Z(\mathrm{Pt}) + Z(\mathrm{Al_2O_3}) = 78 + 50 = \mathbf{128 = Q_\square(3)}$, а катализатор гидроочистки $\mathrm{Co\text{-}MoS_2}$ соединяет $Z(\mathrm{Co}) = \mathbf{27 = d^3}$ со второй сферой $Z(\mathrm{Mo}) = \mathbf{42 = S_2(3)}$!
  \end{itemize}
\item \textbf{Разгадка трёх главных «чёрных ящиков» катализа и пребиотической химии:}
  \begin{itemize}[noitemsep]
  \item \emph{Неуловимое состояние $S_4$ цикла Кока в Фотосистеме~II ($\mathrm{Mn_4CaO_5}$, $2\,\mathrm{H_2O}\to\mathrm{O_2}\uparrow + 4\,\mathrm{H^+} + 4\,e^-$):}
    кубановый кластер $\mathrm{Mn_4CaO_5}$ несёт заряд $4\times 25 + 20 + 40 = \mathbf{160 = Q_\square + K_P/2}$, с двумя субстратными молекулами воды $2\,\mathrm{H_2O}$ ($20$) даёт $\mathbf{180 = 12\times 15}$ (ровно $\mu=15$ квантов на каждый из 12 лучей кубооктаэдра!), а с четырьмя молекулами воды $W_1\dots W_4$ даёт $\mathbf{200 = 25\times 8 = d_{\mathrm{out}}(E)^2|{\triangle}|}$. На вспышках $S_1, S_2, S_3$ ($41k = 41, 82, 123 < 128$) оболочка удерживает кислород ($\mathcal{C}_3 > 0$, причём $S_1$ с $\Phi_{\mathrm{gate}}=0$ является тёмно-стабильным состоянием!). Ровно на \textbf{4-й вспышке $S_4$ ($4\times 41 = 164 > 128$)} связь «кластер --- кислород» становится антисвязывающей ($\mathcal{C}_3(164) = -36 < 0$), отталкивая молекулу $\mathrm{O_2}$, тогда как сам оксо-мостик $\mathrm{O5=O6}$ ($\varphi(t=1) = 106$) сохраняет положительную двойную связь $\mathcal{E}_{\mathrm{bond}} = \mathcal{C}_3(106) = \mathbf{+22 > 0}$ и выбрасывается наружу током $\mathcal{J}_{\mathrm{rxn}} = \mathcal{S}_3(106) = \mathbf{+106 > 0}$!
  \item \emph{«Соединение~Q» метанмонооксигеназы ($\mathrm{CH_4 \to CH_3OH}$ при комнатной температуре $293\text{ К}$, $\hat{n}_E = 0$, без сгорания в $\mathrm{CO_2}$):}
    почему ромб $\mathrm{Fe_2^{IV}(\mu\text{-O})_2}$ ($Z = 2\times 26 + 16 = \mathbf{68 = K_P+4}$) легко рвёт прочную связь $\mathrm{C-H}$ в метане ($439\text{ кДж/моль}$), но не трогает более слабую связь ($402\text{ кДж/моль}$) в получившемся метаноле $\mathrm{CH_3OH}$? Потому что комплекс $[\mathrm{Fe_2O_2 + CH_4}]$ ($68 + 10 = \mathbf{78 = 6\times 13 = Z(\mathrm{Pt})}$) входит в платиновый резонанс, но как только один атом кислорода ($Z=8$) переходит на метан, образующийся метанол $\mathrm{CH_3OH}$ \textbf{замыкает 18-электронную благородную оболочку аргона $Z(\mathrm{CH_3OH}) = 6 + 4 + 8 = \mathbf{18 = |N| + |\square| = Z(\mathrm{Ar})}$} (оставляя $\mathrm{Fe_2O}$ на $Z = 60 = 5\times 12$) и мгновенно отщепляется как фазово-инертный продукт!
  \item \emph{Автокатализ Бутлерова (формозная реакция) и секрет боратной стабилизации D-рибозы РНК ($5\,\mathrm{CH_2O}\to\mathrm{C_5H_{10}O_5}$):}
    формальдегид $\mathrm{CH_2O}$ имеет $Z = \mathbf{16 = 2|{\triangle}|}$, образуя лестницу сахаров $Z(k) = 16k$ ($k=2$ димер-автокатализатор $\mathrm{C_2H_4O_2}$ $Z = \mathbf{32 = K_P/2}$; $k=4$ эритроза $Z = \mathbf{64 = K_P}$; $k=5$ \textbf{D-рибоза РНК} $Z = \mathbf{80 = 2(\Delta\varphi_{\mathrm{disc}}-1)}$; $k=6$ глюкоза $Z = \mathbf{96 = |N||{\triangle}|}$). Присутствие бора ($Z=5=d_{\mathrm{out}}(E)$) останавливает осмоление ровно на $k=5$ рибозе, потому что три водные формы бора точно совпадают с замками носителя: метаборат $\mathrm{BO_2^-}$ имеет $Z = 5+16 = \mathbf{21 = \Delta\varphi_{\mathrm{seam}}(3)}$, борная кислота $\mathrm{B(OH)_3}$ имеет $Z = 5+27 = \mathbf{32 = K_P/2}$, а \textbf{тетрагидроксоборат $\mathrm{B(OH)_4^-}$ имеет $Z = 5 + 4\times 9 = \mathbf{41 = \Delta\varphi_{\mathrm{disc}}(3)}$ (точный замок мёртвой зоны Гейзенберга!)}, замыкая с рибозой ($80$) и $\mathrm{H^+}$ благородную оболочку радона $80 + 5 + 1 = \mathbf{86 = Z(\mathrm{Rn})}$, а с пентамером $\mathrm{HCN}$ --- аденином $\mathrm{C_5H_5N_5}$ ($Z = 5\times 14 = \mathbf{70 = K_P+|\square|}$) --- собирая молекулу АТФ ($Z = 80 + 70 - 10 + 120 = \mathbf{260 = 4\times 5\times 13}$)!
  \end{itemize}
\end{enumerate}

\section{Урок 18. Куперовские пары \texorpdfstring{$2e$}{2e}, нитинол \texorpdfstring{$\mathrm{TiNi}$}{TiNi} и высокотемпературные сверхпроводники}

Почему одиночные электроны ($1e$) на полупериоде Паули $\Delta\varphi = 256$ гасят друг друга ($\mathcal{C}_3(256) = -128$), а спаренные электроны Леона Купера ($2e$, 1956) дают сверхпроводимость без сопротивления?

Подадим удвоенный заряд куперовской пары $(c_0/c)^2 = 2$ в наш ротатор:
\[
  \mathcal{C}_3(2\cdot\Delta\varphi)\Big|_{\Delta\varphi = 128} \;=\; \mathcal{C}_3(256) \;=\; -128 \quad (\text{минимум на четверть-обороте } 128),
\]
но на самом полупериоде Паули $\Delta\varphi = \Phi_{\mathrm{Pauli}} = 256$:
\begin{equation}
  \label{eq:course-cooper-revival}
  \mathcal{C}_3(2\times 256) \;=\; \mathcal{C}_3(512) \;=\; \mathcal{C}_3(0) \;=\; \mathbf{+128},
  \qquad
  \mathcal{S}_3(2\times 256) \;=\; \mathbf{0}!
\end{equation}
На замке Паули $\Delta\varphi = 256$ куперовская пара испытывает \textbf{полное конструктивное возрождение Ааронова--Бома} (квант магнитного потока сверхпроводника $\Phi_0 = h / 2e$)! При этом формульные заряды ключевых функциональных сплавов и сверхпроводников замыкают точные оболочки носителя: сплав с памятью формы нитинол $\mathrm{TiNi}$ имеет $Z(\mathrm{Ti})+Z(\mathrm{Ni}) = 22 + 28 = \mathbf{50}$ (та же вторая высокая магическая оболочка, что у рубина/сапфира $\mathrm{Al_2O_3}$ и перовскита мантии Земли $\mathrm{MgSiO_3}$!), диборид магния $\mathrm{MgB_2}$ имеет $Z = 12 + 2\times 5 = \mathbf{22 = Z(\mathrm{Ti})}$, купратная плоскость $\mathrm{CuO_2}$ имеет $Z = 29 + 16 = \mathbf{45 = 9\times 5}$, а бесконечноелевой ВТСП $\mathrm{CaCuO_2}$ имеет $Z = 20 + 29 + 16 = \mathbf{65 = K_P + 1}$!

